\documentclass[12pt]{article}
\usepackage{amsmath, amssymb}
\usepackage{geometry}
\geometry{margin=1in}

\title{\textbf{Lecture Scribe}}
\author{
Name: Sakshi Patel\\
ID: AU2440206\\
Course: Fundamentals of Probability in Computing\\
}
\date{}

\begin{document}

\maketitle

\section*{�� PART 1: MAIN TOPICS }

From your lecture, here’s the logical flow:
\begin{itemize}
\item Random Vectors (Multivariate Random Variables)
\item Mean Vector
\item Covariance \& Correlation Matrix
\item Linear Transformations $Y = AX + b$
\item Expectation under Transformation
\item Covariance under Transformation
\item Multivariate Gaussian Distribution
\item Special Cases (Zero Correlation vs Independence)
\item Case Study: Smart City Traffic + Weather System
\end{itemize}

\section*{�� PART 2: TEACHING MODE }

\section*{1️⃣ RANDOM VECTORS}

\textbf{Concept}

A random vector is a collection of multiple random variables treated as a single entity.

\[
X =
\begin{bmatrix}
X_1 \\
X_2 \\
\vdots \\
X_n
\end{bmatrix}
\]

\textbf{Why do we need this?}

Because real-world systems are multi-dimensional.

One variable = oversimplification \\
Multiple variables = reality

\textbf{Example}

\[
X =
\begin{bmatrix}
\text{Rainfall} \\
\text{Traffic} \\
\text{Speed}
\end{bmatrix}
\]

\textbf{Intuition}

A state of the system at one moment

\section*{2️⃣ MEAN VECTOR}

\textbf{Definition}

\[
\mu = E[X] =
\begin{bmatrix}
E[X_1] \\
E[X_2] \\
\vdots
\end{bmatrix}
\]

\textbf{Step-by-step reasoning}

Each variable has an average \\
Combine them into a vector \\
That becomes the system’s “center”

\textbf{Example}

\[
\mu =
\begin{bmatrix}
10 \\
50 \\
40
\end{bmatrix}
\]

\section*{3️⃣ COVARIANCE MATRIX}

\textbf{Definition}

\[
\Sigma = E[(X - \mu)(X - \mu)^T]
\]

\textbf{Two variables case}

\[
\text{Cov}(X,Y) = E[(X - \mu_X)(Y - \mu_Y)]
\]

\textbf{Matrix form}

\[
\Sigma =
\begin{bmatrix}
\text{Var}(X_1) & \text{Cov}(X_1, X_2) \\
\text{Cov}(X_2, X_1) & \text{Var}(X_2)
\end{bmatrix}
\]

\textbf{Interpretation}

Covariance matrix = Dependency blueprint of the system

\section*{4️⃣ LINEAR TRANSFORMATION}

\textbf{Definition}

\[
Y = AX + b
\]

\textbf{Example}

\[
Y = 0.7 \cdot \text{Traffic} - 0.3 \cdot \text{Speed}
\]

\section*{5️⃣ EXPECTATION UNDER TRANSFORMATION}

\textbf{Formula}

\[
E[Y] = A\mu + b
\]

\textbf{Derivation}

\[
Y = AX + b
\]

\[
E[Y] = E[AX + b]
\]

\[
= A E[X] + b
\]

\[
= A\mu + b
\]

\section*{6️⃣ COVARIANCE UNDER TRANSFORMATION}

\textbf{Formula}

\[
\Sigma_Y = A \Sigma_X A^T
\]

\textbf{Derivation}

\[
Y = AX + b
\]

\[
Y - \mu_Y = A(X - \mu_X)
\]

\[
(Y - \mu_Y)(Y - \mu_Y)^T
\]

\[
= A(X - \mu_X)(X - \mu_X)^T A^T
\]

\[
\Sigma_Y = A \Sigma_X A^T
\]

\section*{7️⃣ MULTIVARIATE GAUSSIAN}

\textbf{Definition}

A random vector whose joint distribution is Gaussian.

\textbf{Property}

\[
X \sim N(\mu, \Sigma)
\]

\[
Y = AX + b \sim N(A\mu + b, A\Sigma A^T)
\]

\section*{8️⃣ SPECIAL CASES}

Zero covariance $\neq$ independence

Cov = 0 → no linear relation \\
Nonlinear relation can exist

If Gaussian:

Zero covariance ⇒ independence

\section*{9️⃣ CASE STUDY}

\[
X =
\begin{bmatrix}
\text{Rain} \\
\text{Traffic} \\
\text{Speed} \\
\text{Pollution} \\
\text{Temperature}
\end{bmatrix}
\]

\textbf{Transformation}

\[
Y = AX
\]

\section*{�� FINAL SUMMARY}

\begin{itemize}
\item Random vector = multiple variables together
\item Mean = average state
\item Covariance = relationships
\item Linear transform = combine variables
\item Mean transforms linearly
\item Covariance transforms as $A \Sigma A^T$
\item Gaussian = defined by mean + covariance
\item Zero covariance $\neq$ independence
\end{itemize}

\section*{�� FORMULA SHEET}

\[
\mu = E[X]
\]

\[
\Sigma = E[(X - \mu)(X - \mu)^T]
\]

\[
Y = AX + b
\]

\[
E[Y] = A\mu + b
\]

\[
\Sigma_Y = A \Sigma_X A^T
\]

\section*{�� PROBABLE EXAM QUESTIONS}

\begin{itemize}
\item Define random vector with example
\item Derive $E[Y] = A\mu + b$
\item Derive covariance transformation
\item Explain covariance matrix interpretation
\item Difference: covariance vs independence
\item Multivariate Gaussian properties
\item Numerical on linear transformation
\item Case-study based conceptual question
\end{itemize}

\end{document}
